% Proposed receiving source-comparison remark. Root owns insertion and numbering.
\begin{remark}[Exact transports to the cited elliptic and toroidal constructions]
\label{iv:source_transports}
The elliptic quotient and logarithmic-transform construction used here is
the construction of \cite[Section~2 and Remark~8.5]{Engel2026}.
Its periodic toroidal filling is \cite[Proposition~2.8]{Engel2026}.
The following maps fix the scope of that credit.

For the original member
\[
 E_0:\ y^2=x^3+tx+t^2,\qquad G=(0,t),
\]
put
\[
 \tau=\frac4{27t+4},\qquad
 X=\frac92\tau^2x,\qquad
 Y=\frac{27}{2\sqrt2}\tau^3y.
\]
The inverse on the common ordinary open is
\[
 t=\frac{4(1-\tau)}{27\tau},\qquad
 x=\frac{2X}{9\tau^2},\qquad
 y=\frac{2\sqrt2Y}{27\tau^3}.
\]
Substitution gives exactly Engel's displayed equation
\[
 Y^2=X^3-3\tau^3(\tau-1)X+2\tau^4(\tau-1)^2,
\]
and sends $G$ to
$P=(0,\sqrt2\tau^2(1-\tau))$ and $O$ to $O$.
Choose this sign of Engel's primitive generator.
At $\tau=0$, the original minimal coordinates
$x_\infty=t^{-2}x$, $y_\infty=t^{-3}y$ satisfy
\[
 x_\infty=\frac{81X}{8(1-\tau)^2},\qquad
 y_\infty=\frac{729\sqrt2Y}{32(1-\tau)^3}.
\]
These are invertible analytic changes there. The displayed changes are
also units at $\tau=1$. At $\tau=\infty$, Engel's minimal coordinates
are $X/\tau^2=(9/2)x$, $Y/\tau^3=27y/(2\sqrt2)$, and its base
coordinate is $1/\tau=(27/4)(t+4/27)$.
Thus the maps and their inverses identify the minimal regular sources,
the marked point divisors, and the specified singular-fibre germs.
Pulling back an actual root diagram uses the analytic root of the
displayed base unit with its prescribed reference value; its reciprocal
is the inverse root comparison. The finite character is retained.
Engel's linearization is a section of
$\mathcal Hom(t_{6P}^*M,M)$; Definition~2.3 explicitly makes its
bundle automorphism cover $t_{6P}^{-1}$. Under this map that source
translation is exactly $-6G$.

There is also an explicit relation among the four regular point-line
constructions. Write $M_a=\mathcal O([aG]-[O])$, $a\in\{1,2,3,6\}$,
$b=6/a$, and retain the actual Miller function
\[
 \operatorname{div}(f_a)=a[G]-[aG]-(a-1)[O].
\]
Multiplication by $f_a$ is the line isomorphism
$M_1^{\otimes a}\longrightarrow M_a$; consequently
\[
 \mathcal F_a(p,v)=(p,f_a(p)v^a)
\]
is its degree-$a$ map on nonzero lines. The stated divisor equality
proves the map at zeros and poles in the point Cartier modules.
Put $T(p)=p-G$ and
\[
 B_{a,1}(p)=f_2(p)^a\frac{f_a(Tp)}{f_a(p)},\qquad
 B_{a,b}(p)=\prod_{j=0}^{b-1}B_{a,1}(T^jp).
\]
In the common Miller normalization let
\[
 \mathcal B_a(p,v)=
 \left(p-bG,\lambda\frac{q}{t^3}B_{a,b}(p)v\right),
 \qquad q=t+4/27.
\]
The chronological product telescopes and gives the literal identity
\[
 \mathcal F_a\mathcal B_1=\mathcal B_a^{\circ a}\mathcal F_a.
\]
Indeed both line coefficients are
\[
 f_a(p-6G)
 \left(\lambda\frac q{t^3}
       \prod_{j=0}^{5}f_2(T^jp)\right)^a.
\]
For specified sections with nonzero reference constants $C_1,C_a$,
use the exact scalar transport $\lambda_a C_a=\lambda_1 C_1$.
The induced map on every regular compact two-torus has degree $a^2$.
Indeed, the displayed point-Cartier map has degree $a$ on the nonzero
line. Equivariance gives its degree-$a$ quotient map to
$M_a^\times/\langle\mathcal B_a^{\circ a}\rangle$; the latter maps
with degree $a$ to $M_a^\times/\langle\mathcal B_a\rangle$, because
the stated suspension subgroup has index $a$ in its infinite cyclic
deck group. This factorization uses the entire Miller coefficient and
fixes the degree without choosing logarithmic elliptic lifts. It
specifies the regular-locus comparison. Each statement's own root and
fan maps provide its claimed extension across the completed pieces.

For the other rational elliptic sources, the application of Engel's
Remark~8.5 is equally explicit. Its $S,C,M,t$ are the displayed
minimal regular source, its original base, the specified Cartier point
line (including any stated vertical component divisor), and the
translation by $bG$. The Hom section defines the actual bundle map
over $-bG$, as in Definition~2.3. For the weighted source in
\cref{night:P3F.21}, these data are $M=\mathcal O(7G-O)$ and
$bG=2G$; the displayed source intersection calculation gives the
required Hom bundle $\pi^*\mathcal O(1)$. Its period generators
$h,g$ are changed to $H=h^{-1}g^{-2}$ and $K=h^4g^7$ by the
matrix
\[
 \begin{pmatrix}-1&4\\-2&7\end{pmatrix}\in SL_2(\mathbb Z).
\]
This is the actual quotient-lattice map used at that cusp.

For the reference cusp, \cite[Proposition~9.24 and
Remarks~9.25--9.26, pp.~82--83]{AC26} already computes the fixed
rational curve, its weights and every finite-subgroup fixed set.
Write $T_+(u,v)=\{(0,u,u+v),(0,v,u+v)\}+\mathbb Z^2$
for the positive-diagonal convention of \cref{night:P3F.35}.
The AC26 triangles $(0,e_1,e_2)$ and $(e_1,e_2,e_1+e_2)$ are exactly
$T_+(e_1,-e_2)$ after translating its two representatives by $e_2$;
the physical $\delta=e_2$ is $-w$ in this basis. Thus the adapter
is the integral matrix $\operatorname{diag}(1,-1)$, together with
that explicit lattice translation.
On a general basis fan the same map is
$\mathbb Z^2\to\mathbb Z^2/\mathbb Ze$ for an edge $e$:
the image of the physical cocharacter is $\det(e,\delta)$.
The basis change in \cref{night:P3F.35} carries that literal quotient
map to the three displayed edge orders. Its common blowup and two
ruling contractions carry the integral topology to every fan.
The final cap-freeness deduction uses
\cite[Theorem~4.4, p.~146; coefficient convention p.~139]{Floyd1952}
with those actual prime fixed sets.
\end{remark}
